{
    \small
    \renewcommand{\arraystretch}{1.3}
    \begin{longtable}{|>{\raggedright\arraybackslash}p{1.8cm}|>{\raggedright\arraybackslash}p{2.3cm}|>{\raggedright\arraybackslash}p{6.6cm}|>{\raggedright\arraybackslash}p{3.1cm}|}
    \caption{Requisiti di qualità}\label{tab:req-qualita}\\
    \hline
    \rowcolor[gray]{0.88} \textbf{Codice} & \textbf{Importanza} & \textbf{Descrizione} & \textbf{Fonti} \\ \hline
    \endfirsthead
    \hline
    \rowcolor[gray]{0.88} \textbf{Codice} & \textbf{Importanza} & \textbf{Descrizione} & \textbf{Fonti} \\ \hline
    \endhead
    \mbox{R1-Q-O} & Obbligatorio & Il sistema deve essere corredato di \textit{test}\textsubscript{G} di unità e di integrazione che verifichino l'affidabilità dell'intera pipeline end-to-end ed evitino regressioni. & \mbox{Capitolato}, \mbox{LG-5.1} \\ \hline
    \mbox{R2-Q-O} & Obbligatorio & L'architettura deve essere modulare, così da consentire l'integrazione agevole di nuovi canali di ingresso o nuove funzionalità in futuro. & \mbox{Capitolato}, \mbox{LG-2.1} \\ \hline
    \mbox{R3-Q-O} & Obbligatorio & L'architettura deve essere scalabile e in grado di sostenere l'elaborazione di elevati volumi di richieste e di dati sensibili. & \mbox{Capitolato}, \mbox{VE-1.9} \\ \hline
    \mbox{R4-Q-O} & Obbligatorio & Nell'elaborazione delle richieste l'affidabilità del processo ha priorità assoluta sulla velocità di esecuzione. & \mbox{VE-1.9} \\ \hline
    \mbox{R5-Q-O} & Obbligatorio & I parametri applicativi, i canali e le regole di notifica devono essere gestiti tramite configurazioni esternalizzate (JSON o YAML). & \mbox{Capitolato} \\ \hline
    \mbox{R6-Q-O} & Obbligatorio & Il modello LLM deve poter essere sostituito a caldo con un altro modello, senza modificare la logica applicativa. & \mbox{VE-1.8}, \mbox{LG-2.2} \\ \hline
    \mbox{R7-Q-O} & Obbligatorio & Il trattamento dei dati deve essere conforme al \textit{GDPR}\textsubscript{G}, adottando i principi di security by default e security by design. & \mbox{LG-3} \\ \hline
    \mbox{R8-Q-O} & Obbligatorio & Il progetto deve essere corredato dalla documentazione prevista: diagrammi \textit{UML}\textsubscript{G} dei casi d'uso, schema della base dati o delle API richiamate, casi di test con relativa reportistica e lista dei bug risolti. & \mbox{Capitolato} \\ \hline
    \mbox{R9-Q-O} & Obbligatorio & Il codice sorgente deve essere gestito tramite un sistema di versionamento, quale GitHub o Bitbucket. & \mbox{Capitolato} \\ \hline
    \mbox{R10-Q-O} & Obbligatorio & I test devono utilizzare il \textit{dataset}\textsubscript{G} \textit{MOCK}\textsubscript{G} fornito dalla proponente (e-mail, PDF, risposte JSON). & \mbox{VE-1.1}, \mbox{LG-5.2} \\ \hline
    \mbox{R11-Q-O} & Obbligatorio & Il gruppo deve produrre registrazioni audio di prova per WhatsApp, simulando contesti reali d'uso in mobilità con rumori di fondo. & \mbox{LG-5.2} \\ \hline
    \mbox{R12-Q-D} & Desiderabile & La trascrizione dei messaggi vocali dovrebbe risultare robusta ai rumori di fondo tipici di contesti d'uso in mobilità. & \mbox{LG-5.2} \\ \hline
    \end{longtable}
}